\section{Marco teórico}

\subsection{Simulador trifásico balanceado}
El STB puede modelarse mediante tres fuentes senoidales de igual amplitud y frecuencia, separadas angularmente \SI{120}{\degree}. Tomando la fase $a$ como referencia:
\begin{align}
	v_a(t)&=V_m\sin(\omega t), \\
	v_b(t)&=V_m\sin(\omega t-120^{\circ}), \\
	v_c(t)&=V_m\sin(\omega t+120^{\circ}).
\end{align}

En forma fasorial, los voltajes de fase se representan como:
\begin{align}
	\mathbf{V}_{an}&=V_F\angle 0^{\circ}, \\
	\mathbf{V}_{bn}&=V_F\angle -120^{\circ}, \\
	\mathbf{V}_{cn}&=V_F\angle 120^{\circ}.
\end{align}

Para esta práctica se toma como referencia el valor indicado en el manual del STB: $V_{pp}=\SI{14.4}{V}$ y $f=\SI{925.93}{Hz}$. Así:

\begin{align}
	V_m&=\frac{V_{pp}}{2}=\SI{7.2}{V}, \\
	V_F&=V_{\rms}=\frac{V_m}{\sqrt{2}}=\SI{5.091}{V}, \\
	T&=\frac{1}{f}=\SI{1.080}{ms}, \\
	\Delta t_{120^{\circ}}&=\frac{T}{3}=\SI{0.360}{ms}.
\end{align}

\subsection{Voltajes de línea}
Los voltajes de línea son diferencias fasoriales entre dos fases:
\begin{align}
	\mathbf{V}_{ab}&=\mathbf{V}_{an}-\mathbf{V}_{bn}, \\
	\mathbf{V}_{bc}&=\mathbf{V}_{bn}-\mathbf{V}_{cn}, \\
	\mathbf{V}_{ca}&=\mathbf{V}_{cn}-\mathbf{V}_{an}.
\end{align}

Para un sistema balanceado:
\begin{equation}
	V_L=\sqrt{3}V_F.
\end{equation}

Usando $V_F=\SI{5.091}{V_{RMS}}$:
\begin{equation}
	V_L=\sqrt{3}(\SI{5.091}{V})=\SI{8.818}{V_{RMS}}.
\end{equation}

\subsection{Configuración estrella}
En una conexión estrella, cada carga se conecta entre una línea y un punto común $n'$. Si el neutro de la fuente se conecta al punto común, los voltajes de fase quedan definidos directamente por el STB. Para una carga balanceada resistiva:
\begin{equation}
	I_F=I_L=\frac{V_F}{R}.
\end{equation}

Además, la corriente de neutro se determina por suma fasorial:
\begin{equation}
	\mathbf{I}_N=\mathbf{I}_a+\mathbf{I}_b+\mathbf{I}_c.
\end{equation}

En un sistema balanceado ideal, $\mathbf{I}_N=0$.

\subsection{Configuración delta}
En una conexión delta, cada impedancia se conecta entre dos líneas. Por ello, el voltaje de fase de cada carga es igual al voltaje de línea:
\begin{equation}
	V_F^{(\Delta)}=V_L.
\end{equation}

Para una delta balanceada:
\begin{equation}
	I_L=\sqrt{3}I_F.
\end{equation}

\subsection{Impedancia capacitiva y triángulo de impedancias}
Para el capacitor de la práctica:
\begin{equation}
	X_C=\frac{1}{2\pi f C}.
\end{equation}

Con $f=\SI{925.93}{Hz}$ y $C=\SI{0.22}{\micro F}$:
\begin{equation}
	X_C=\frac{1}{2\pi(925.93)(0.22\times 10^{-6})}=\SI{781.30}{\ohm}.
\end{equation}

Si se emplea una rama serie $R$-$C$ con $R=\SI{1}{\kilo\ohm}$:
\begin{align}
	Z_{RC}&=R-jX_C=1000-j781.30\;\Omega, \\
	|Z_{RC}|&=\SI{1269.03}{\ohm}, \\
	\theta&=-\tan^{-1}\left(\frac{781.30}{1000}\right)=-38.00^{\circ}.
\end{align}

\begin{figure}[H]
	\centering
	\shorthandoff{<>}
	\begin{tikzpicture}[scale=1.0,>=Latex]
		\draw[->] (0,0) -- (5,0) node[right]{$R=1000\,\Omega$};
		\draw[->] (5,0) -- (5,-3.9) node[midway,right]{$-X_C=-781.3\,\Omega$};
		\draw[->,thick] (0,0) -- (5,-3.9) node[midway,below left]{$|Z|=1269.03\,\Omega$};
		\draw (1.1,0) arc[start angle=0,end angle=-37.95,radius=1.1];
		\node at (1.6,-0.45) {$\theta=-38^{\circ}$};
	\end{tikzpicture}
	\caption{Triángulo de impedancias para una rama serie $R$-$C$ usando el capacitor de \SI{0.22}{\micro F}.}
\end{figure}